\chapter{Project Overview}

\section{Engineering problem}
State the practical engineering problem in terms a hiring manager, customer, or
technical reviewer can evaluate. Identify the operating environment, performance
objective, constraints, and consequences of failure. Replace generic phrases such as
``optimize the system'' with measurable engineering targets.

\section{System context}
Figure~\ref{fig:portfolio-workflow} shows the recommended evidence flow for a public
engineering project. The report begins with requirements, proceeds through analytical
and numerical work, and ends with controlled public claims.

\begin{figure}[H]
\centering
\begin{tikzpicture}[
  node distance=7mm and 9mm,
  box/.style={draw,rounded corners,minimum width=2.65cm,minimum height=1.05cm,
              align=center,fill=blue!5,font=\small},
  arr/.style={-{Latex[length=2.4mm]},thick}
]
\node[box] (req) {Requirements\\and constraints};
\node[box,right=of req] (ana) {Analytical\\design};
\node[box,right=of ana] (model) {Model or\\prototype};
\node[box,below=of model] (verify) {Verification\\and validation};
\node[box,left=of verify] (judge) {Engineering\\interpretation};
\node[box,left=of judge] (claim) {Controlled\\public claims};
\draw[arr] (req) -- (ana);
\draw[arr] (ana) -- (model);
\draw[arr] (model) -- (verify);
\draw[arr] (verify) -- (judge);
\draw[arr] (judge) -- (claim);
\end{tikzpicture}
\caption{Recommended evidence flow for an engineering portfolio project.}
\label{fig:portfolio-workflow}
\end{figure}

\section{Portfolio value}
Explain what this project demonstrates beyond software operation. Examples include
requirements translation, model credibility, trade-space reasoning, numerical
reliability, root-cause analysis, design iteration, verification discipline, and
technical communication.
